\section{Discussion}
\label{sec:discussion}

\subsection{Key Findings}

\textbf{Finding 1: Venue-string classification is not a reliable method for interdisciplinary alignment corpora.}

The 12\% coverage finding is a blocking problem, not a minor detail.
A bibliometric study that classifies only 12\% of its corpus papers cannot construct a valid $2\times2$ contingency table.
The failure occurs silently: the classification function returns \texttt{None} for unmatched papers, and downstream analysis omits them without warning.
Researchers must use FoS-primary classification or normalize venue strings to full proceedings name variants before using S2AG data for interdisciplinary corpus classification.

\textbf{Finding 2: The directional signal is visible even in the proxy corpus.}

A ratio of 0.107 at $N=9$ cannot be reported as a confirmed statistical finding, and we are explicit about this.
But even in 33 papers, the asymmetry predicted by H-CitAsym is visible in the raw edge counts, and no falsifier was triggered.
The direction has not been refuted.

\textbf{Finding 3: Public reading lists are not bibliometric dataset proxies.}

The 12.5\% yield from GitHub-linked papers to S2AG-resolved papers is a practical finding for researchers in adjacent areas.
Programmatic S2AG search rather than reading list parsing is required to reconstruct systematic review corpora at bibliometric scale.

\subsection{Limitations}

\textbf{Limitation 1: All statistical predictions (P1, P2, P3) are INCONCLUSIVE --- corpus too small for chi-squared testing.}

Coverage = 67.3\% and cross-group edges = 9 both fall below pre-registered thresholds.
The h-e1-v2 protocol specifies a concrete executable augmentation step: use S2AG \texttt{/paper/arXiv:2406.09264/citations} to discover papers citing \citet{Shen2024Position} and expand the corpus programmatically.

\textbf{Limitation 2: Study operates on 33-paper proxy corpus, not the intended 400-paper systematic review corpus.}

The corpus proxy mismatch is itself a finding (Section~\ref{sec:results}).
Future researchers now have a documented account of why GitHub reading lists fail as bibliometric proxies.

\textbf{Limitation 3: Sensitivity analysis across 3 schemes is not executable at Scheme 1/2 coverage of 12\%.}

The fix is mechanical: extend venue substring sets to full proceedings name variants.
This fix is documented in \texttt{config.py} but not applied during h-e1 execution.

\textbf{Limitation 4: Findings apply specifically to the alignment community as operationalized by \citet{Shen2024Position}.}

Field-level generalization requires a matched baseline.
Phase 5 comparison is deferred by pipeline configuration.

\textbf{Limitation 5: Temporal dimension unverified.}

Pre-/post-2022 cohort analysis is not meaningful at $N=33$.
The full corpus (h-e1-v2) will enable this analysis using the \texttt{year} field already retrieved in batch resolution.

\subsection{Broader Impact}

This study contributes a validated S2AG bibliometric pipeline and FoS-primary classification method applicable to any interdisciplinary corpus indexed by S2AG.
All code and cached API responses are released as supplementary material.
Readers should not cite the ratio 0.107 as confirmed evidence of asymmetry without acknowledging the corpus scale limitation.
There is no foreseeable negative impact from measuring citation behavior in a public academic corpus.
